\newif\ifglobal

%%%%%%%%%%%%%%%%%%%%%%%
%%% Compile to view %%%
%%%%%%%%%%%%%%%%%%%%%%%

%%%%%%%%%%%%%%%%%%%%%%%%%%%%%%%%%%%%%%%%%%%%%%%%%%%
%%% Readme:                                     %%%
%%% https://www.overleaf.com/read/bwdjvfjksvjy/ %%%
%%%%%%%%%%%%%%%%%%%%%%%%%%%%%%%%%%%%%%%%%%%%%%%%%%%

% >>> M?: marks a module setting number or important notice

% >>> M4: Set {./relative-path} to external files for this module <<<
\def \modulefiles {./24-LPMAILBOX}
% To add an external file, use:
% (Replace '---' with actual filename)
% '\input{\modulefiles/tables/---.tex}' for tables
% '\includegraphics{\modulefiles/figures/---}' for figures
% '\subfile{\modulefiles/---}' for register files


% >>> M1: This module is in global project? <<<
% 'YES' - uncomment; 'NO' - comment out
\globaltrue

\ifglobal
    \documentclass[main\_\_CN.tex]{subfiles}

\else
    % Font "FandolSong-Regular" used by ctex does not have GSUB table.
    % It causes warnings that do not affect compilation and can be ignored.
    \PassOptionsToPackage{quiet}{fontspec}

    \documentclass[openany, 10pt]{book}

    % >>> M2: In ./00-shared/preamble-trm-module, also find
    %         settings for visibility of todo notes, watermark <<<
    \usepackage{preamble-trm-module}

    % Enable Chinese typesetting
    \usepackage{ctex}

    % >>> M3: Temporary labels in Chinese
    %         you can add your own labels there <<<
    \myexternaldocument{./00-shared/temp-labels-cn}

\fi %\ifglobal


\setcounter{secnumdepth}{4}


% >>> M5: If this module needs any updates in
% './00-shared/preamble-trm-module.sty'
% (include additional package, etc.), contact your module PM
% or create the file './00-shared/changelog.tex' make the
% necessary updates yourself and document them in this file <<<

% Make text left-aligned and allow latex to break pages naturally, comment out for full justification
\raggedright
\raggedbottom


\begin{document}


% Sets language into which pre-defined bits of text are translated
\selectlanguage{chinese}
\setchineseformatting


% Do not delete this block
\ifglobal
% Add the following front matter if
% this module is outside of global project
\else
    \InputIfFileExists{readme}{}{}
    \clearpage
    \listoftodos
    \listofdonetodos
    \tableofcontents
    %\listoftables
    %\listoffigures
\fi


%%%%%%%%%%%%%%%%%%%%%%%%%
%%% TRM Module Begins %%%
%%%%%%%%%%%%%%%%%%%%%%%%%

\hypertarget{lpmailbox}{}
\chapter{LP 信箱控制器}
\label{mod:lpmailbox}


\section{概述}

\chipname{} 中包含一个 LP 信箱控制器模块，通过硬件机制实现 LP CPU 和 HP CPU0/1 之间高效的核间通信。LP 信箱控制器模块中包含 16 个 32 位信息寄存器可供 LP CPU 和 HP CPU0/1 存储并传递信息，并通过中断机制实现 LP CPU 和 HP CPU0/1 之间的核间通信。

\section{特性}
\chipname{} 的 LP 信箱控制器模块具有以下特性：
\begin{itemize}
    \item 支持多达 16 个 32 位信息寄存器用于核间通信
    \item 支持 LP CPU 外部中断信号 \hyperref[tab:lpmailbox-int-sources]{MB\_LP\_INTR}
    \item 支持 HP CPU\regindex{x} (\regindex{x}: 0$\sim$1) 外部中断信号 \hyperref[tab:lpmailbox-int-sources]{MB\_HP\_INTR}
\end{itemize}

\section{功能描述}
\begin{figure}[H]
    \centering
    \includegraphics[width=0.8\linewidth]{\modulefiles/figures/LP_MB}
    \caption{LP 信箱控制器系统框图}
    \label{fig:lpmbsysdigram}
\end{figure}

\subsection{信息寄存器}
LP 信箱控制器中包含了 \hyperref[regdesc:MBMESSAGE0REG]{MB\_MESSAGE\_0\_REG}$\sim$\hyperref[regdesc:MBMESSAGE15REG]{MB\_MESSAGE\_15\_REG} 共计 16 个 32 位信息寄存器，用于缓存核间通信信息。LP CPU 和 HP CPU0/1 均可读写这些信息寄存器以传递或获取信息。

读写信息寄存器通过 APB 总线来完成，使用 LP\_PERI\_CLK 时钟。该时钟为 LP 信箱控制器使用的唯一时钟，更多时钟信息，见章节 \ref{mod:resclk} \textit{\nameref{mod:resclk}}。

\subsection{中断}
\label{sec:lpmailbox-interrupt}

\chipname{} 的 LP 信箱控制器可以生成如下外部中断信号：
    \begin{itemize}
        \item MB\_LP\_INTR
        \item MB\_HP\_INTR
    \end{itemize}
MB\_LP\_INTR 直接映射至 LP CPU（见章节 \ref{mod:ulp} \textit{\nameref{mod:ulp}}），MB\_HP\_INTR 可通过\hyperref[mod:intmtrx]{中断矩阵}映射至 HP CPU0/1。
这些中断信号由 LP 信箱控制器模块的内部中断源生成。表 \ref{tab:lpmailbox-int-sources} 列出了 LP 信箱控制器的内部中断源、中断使能、中断触发条件、中断清除、以及生成的中断信号。当满足中断使能条件的时候，中断触发条件满足后，内部中断源才会产生。当中断清除条件满足的时候，内部中断源会被清除。

\begin{longtable}{ | L{2.3cm} | L{3.8cm} | L{3.9cm} | L{4.2cm} | L{2.2cm} | }
\caption{LP 信箱控制器的内部中断源} \label{tab:lpmailbox-int-sources}
\\ \hline
\rowcolor{lightgray}
\textbf{内部中断源} & \textbf{中断使能条件}  & \textbf{中断触发条件}  & \textbf{中断清除条件}  &  \textbf{中断信号}\\ \hline
\label{int:MBHPmINT} MB\_HP\_\regindex{i}\_INT & \hyperref[fielddesc:MBHP0INTENA]{MB\_HP\_\regindex{i}\_INT\_ENA} 为 1 &  写 MB\_MESSAGE\_\regindex{i}\_REG & 向 \hyperref[fielddesc:MBHP0INTCLR]{MB\_HP\_\regindex{i}\_INT\_CLR} 写 1  & MB\_HP\_INTR  \\ \hline
\label{int:MBLPnINT} MB\_LP\_\regindex{i}\_INT & \hyperref[fielddesc:MBLP0INTENA]{MB\_LP\_\regindex{i}\_INT\_ENA} 为 1 &  写 MB\_MESSAGE\_\regindex{i}\_REG & 向 \hyperref[fielddesc:MBLP0INTCLR]{MB\_LP\_\regindex{i}\_INT\_CLR} 写 1  & MB\_LP\_INTR  \\ \hline
\end{longtable}
\begin{tiplisting}
\vspace{-2em}
\begin{itemize}
    \item \regindex{i} 取值范围是 0$\sim$15。
    \item LP 信箱控制器并不识别是哪一个 CPU 在访问它，只要满足表 \ref{tab:lpmailbox-int-sources} 的触发条件，中断就会产生或被清除。
    \item\hyperref[tab:lpmailbox-int-sources]{MB\_HP\_INTR} 同样可作为 LP CPU 的外部中断源，通过 LP 中断映射模块映射（LP\_INTR\_MAP，见图 \ref{fig:lpmbsysdigram}）至 LP CPU。有关 LP\_INTR\_MAP 的更多信息，见章节 \ref{mod:ulp} \textit{\nameref{mod:ulp}}。
    \item\hyperref[tab:lpmailbox-int-sources]{MB\_LP\_INTR} 同样可以作为 HP CPU\regindex{x} 的外部中断源，可通过\hyperref[mod:intmtrx]{中断矩阵}将其映射至 HP CPU\regindex{x}。
    \item 当 LP 信箱控制器需要为 HP CPU0 和 HP CPU1 分配不同的中断源时，可通过\hyperref[mod:intmtrx]{中断矩阵}将 \hyperref[tab:lpmailbox-int-sources]{MB\_HP\_INTR} 和 \hyperref[tab:lpmailbox-int-sources]{MB\_LP\_INTR} 分别映射至 HP CPU0 和 HP CPU1。
    \end{itemize}
\end{tiplisting}

\subsection{核间通信}
当 LP CPU 与 HP CPU\regindex{x} 之间需要通信时，可以根据需求利用 LP 信箱控制器中部分或者全部信息寄存器来分享信息，并通过相应的中断来同步信息。在通信期间，需要软件保证仅发送方可写 LP 信箱控制器中的信息寄存器，防止信息寄存器中的数据被篡改。
\begin{itemize}
    \item 以 HP CPU0 通过 MB\_MESSAGE\_\regindex{i}\_REG 向 LP CPU 发起通信为例：
    \begin{itemize}
        \item 配置 LP 中断映射模块（LP\_INTR\_MAP，见图 \ref{fig:lpmbsysdigram}），将 MB\_LP\_\regindex{i}\_INT 中断映射至 LP CPU。
        \item 置位 \hyperref[fielddesc:MBLP0INTENA]{MB\_LP\_\regindex{i}\_INT\_ENA} 使能 MB\_LP\_\regindex{i}\_INT 中断。
        \item HP CPU0 写 \hyperref[regdesc:MBMESSAGE0REG]{MB\_MESSAGE\_\regindex{i}\_REG}，缓存共享信息。
        \item LP CPU 感知到中断发生，进入中断处理程序。在中断处理程序中检测 \hyperref[regdesc:MBLPINTSTREG]{MB\_LP\_INT\_ST\_REG} 中断状态寄存器发现 \hyperref[fielddesc:MBLP0INTST]{MB\_LP\_\regindex{i}\_INT\_ST} 状态位置位，判定发生了 MB\_LP\_\regindex{i}\_INT 中断。
        \item LP CPU 读取 \hyperref[regdesc:MBMESSAGE0REG]{MB\_MESSAGE\_\regindex{i}\_REG} 获取共享信息。
        \item LP CPU 写 1 到 \hyperref[regdesc:MBLPINTCLRREG]{MB\_LP\_INT\_CLR\_REG} 寄存器的 {MB\_LP\_\regindex{i}\_INT\_CLR} 比特位，会触发中断清除动作。
        \item LP CPU 读取 \hyperref[regdesc:MBLPINTSTREG]{MB\_LP\_INT\_ST\_REG} 中断状态寄存器，直到其 \hyperref[fielddesc:MBLP0INTST]{MB\_LP\_\regindex{i}\_ST} 状态位清零。然后 LP CPU 会退出中断处理程序。
    \end{itemize}
    \item 以 LP CPU 向 HP CPU0 发起通信为例：
    \begin{itemize}
        \item 配置\hyperref[mod:intmtrx]{中断矩阵}，将 MB\_HP\_INTR 映射至 HP CPU0（见章节 \ref{mod:intmtrx} \textit{\nameref{mod:intmtrx}}）。
        \item 置位 \hyperref[fielddesc:MBHP0INTENA]{MB\_HP\_\regindex{i}\_INT\_ENA} 使能 MB\_HP\_\regindex{i}\_INT 中断。
        \item LP CPU 写 \hyperref[regdesc:MBMESSAGE0REG]{MB\_MESSAGE\_\regindex{i}\_REG}，缓存共享信息。
        \item HP CPU0 感知到中断发生，进入中断处理程序。在中断处理程序中检测 \hyperref[regdesc:MBHPINTSTREG]{MB\_HP\_INT\_ST\_REG} 中断状态寄存器，发现 \hyperref[fielddesc:MBHP0INTST]{MB\_HP\_\regindex{i}\_INT\_ST} 状态位置位，判定发生了 MB\_HP\_\regindex{i}\_INT 中断。
        \item HP CPU0 读取 \hyperref[regdesc:MBMESSAGE0REG]{MB\_MESSAGE\_\regindex{i}\_REG} 获取共享信息。
        \item HP CPU0 写 1 到 \hyperref[regdesc:MBHPINTCLRREG]{MB\_HP\_INT\_CLR\_REG} 寄存器的 {MB\_HP\_\regindex{i}\_INT\_CLR} 比特位，会触发中断清除动作。
        \item HP CPU0 读取 \hyperref[regdesc:MBHPINTSTREG]{MB\_HP\_INT\_ST\_REG} 中断状态寄存器，直到其 \hyperref[fielddesc:MBHP0INTST]{MB\_HP\_\regindex{i}\_ST} 状态位清零。然后 HP CPU0 会退出中断处理程序。
    \end{itemize}
\end{itemize}


%%%%%%%%%%%%%%%%%%%%%%%%
%%% Register Summary %%%
%%%%%%%%%%%%%%%%%%%%%%%%

% >>> Update the sample text below and insert
%     the info on register summary into the subfile <<<

\hypertarget{lpmailbox-reg-summ}{}
\section{寄存器列表}
\label{sec:lpmailbox-reg-summ}

本小节的所有地址均为相对于 LP 信箱控制器基地址的地址偏移量（相对地址），具体基地址请见章节 \ref{mod:sysmem} \textit{\nameref{mod:sysmem}} 中的表 \ref{tab:sysmem-base-address}。

请查看章节 \hyperref[glossary-access-types]{\textit{寄存器的访问类型}}，了解“访问”列缩写的含义。

\subfile{\modulefiles/24-LPMAILBOX-reg-summ__CN}

%%%%%%%%%%%%%%%%%
%%% Registers %%%
%%%%%%%%%%%%%%%%%

% >>> Update the sample text below and insert
%      the info on registers into the subfile <<<

\newpage
\section{寄存器}

本小节的所有地址均为相对于 LP 信箱控制器基地址的地址偏移量（相对地址），具体基地址请见章节 \ref{mod:sysmem} \textit{\nameref{mod:sysmem}} 中的表 \ref{tab:sysmem-base-address}。

关于保留 (reserved) 域的处理，请查看章节 \hyperref[programming-reserved-field]{如何配置寄存器的保留域}。

\subfile{\modulefiles/24-LPMAILBOX-reg__CN}

\end{document}
